\documentclass[11pt,a4paper]{article}
\usepackage[T1]{fontenc}
\usepackage[utf8]{inputenc}
\usepackage[french]{babel}
\usepackage[a4paper,margin=2cm,bottom=2.5cm]{geometry}
\usepackage{lmodern}
\usepackage{amsmath,amssymb}
\usepackage{array,booktabs}
\usepackage{enumitem}
\usepackage{multicol}
\usepackage[most]{tcolorbox}
\usepackage{xcolor}
\usepackage{titlesec}
\usepackage{fancyhdr}
\usepackage{needspace}
\usepackage{tikz}
\usetikzlibrary{decorations.pathmorphing}

\definecolor{bleu}{RGB}{29,95,168}
\definecolor{orange}{RGB}{190,95,10}
\definecolor{vert}{RGB}{35,120,75}
\definecolor{rouge}{RGB}{180,40,40}

\setlength{\parindent}{0pt}
\setlength{\parskip}{4pt}
\setlist{nosep,leftmargin=*}
\titleformat{\section}{\Large\bfseries\color{bleu}}{\thesection.}{0.5em}{}
\pagestyle{fancy}\fancyhf{}
\renewcommand{\headrulewidth}{0pt}
\renewcommand{\footrulewidth}{0.4pt}
\fancyfoot[L]{\footnotesize Physique-chimie 1re spécialité --- Constance --- Ondes mécaniques}
\fancyfoot[R]{\footnotesize\thepage}

\newtcolorbox{cours}[1]{colback=bleu!5,colframe=bleu,fonttitle=\bfseries,title={#1},boxrule=0.6pt,arc=2pt,left=6pt,right=6pt}
\newtcolorbox{piege}{colback=rouge!5,colframe=rouge,fonttitle=\bfseries,title=Pièges classiques,boxrule=0.6pt,arc=2pt,left=6pt,right=6pt}
\newcommand{\etoiles}[1]{\ifcase#1\or$\bigstar$\or$\bigstar\bigstar$\or$\bigstar\bigstar\bigstar$\fi}
\newcounter{exo}
\newcommand{\exo}[2]{\par\needspace{8\baselineskip}\medskip\stepcounter{exo}{\bfseries\color{bleu}Exercice \theexo\ --- #2}\hfill{\color{orange}\etoiles{#1}}\par\nopagebreak}
\newcounter{res}
\newcommand{\res}[1]{\par\needspace{4\baselineskip}\smallskip\stepcounter{res}{\bfseries\color{vert}Exercice \theres\ --- #1}\par\nopagebreak}
\newcommand{\ms}{\text{m}\cdot\text{s}^{-1}}
\renewcommand{\arraystretch}{1.3}

\begin{document}

\begin{center}
  {\LARGE\bfseries\color{bleu} Ondes mécaniques}\\[4pt]
  {\small Première, spécialité physique-chimie}\\[2pt]
  {\small L'essentiel du cours, puis des exercices de difficulté croissante :
  {\color{orange}$\bigstar$} application, {\color{orange}$\bigstar\bigstar$} raisonnement, {\color{orange}$\bigstar\bigstar\bigstar$} problème.
  Les résultats sont à la fin.}
\end{center}

\section{L'essentiel du cours}

\begin{cours}{Onde mécanique progressive}
Une \textbf{onde mécanique progressive} est la propagation d'une \textbf{perturbation} dans un \textbf{milieu matériel}, avec transport d'\textbf{énergie} mais \textbf{sans transport de matière}.
\begin{itemize}
  \item Exemples : la ola (les spectateurs restent à leur place), une corde, un ressort, les vagues, le son, les ondes sismiques.
  \item Une onde mécanique \textbf{ne se propage pas dans le vide} (il faut un milieu : solide, liquide ou gaz). Les ondes électromagnétiques (lumière, radio\dots), elles, se propagent aussi dans le vide.
\end{itemize}
\end{cours}

\begin{cours}{Transversale ou longitudinale ?}
\begin{minipage}{0.48\linewidth}
\textbf{Transversale} : la perturbation est \textbf{perpendiculaire} à la direction de propagation.\\
Ex. : corde, ola, surface de l'eau.

\centering
\begin{tikzpicture}[scale=0.7]
  \draw[thick] (0,0) -- (2,0) .. controls (2.5,0) and (2.5,1) .. (3,1) .. controls (3.5,1) and (3.5,0) .. (4,0) -- (6,0);
  \draw[->,orange,thick] (3,1.2) -- (3,1.9) node[right] {\small perturbation};
  \draw[->,bleu,thick] (4.3,0.6) -- (5.8,0.6) node[above] {\small propagation};
\end{tikzpicture}
\end{minipage}\hfill
\begin{minipage}{0.48\linewidth}
\textbf{Longitudinale} : la perturbation est \textbf{parallèle} à la direction de propagation.\\
Ex. : ressort comprimé, son.

\centering
\begin{tikzpicture}[scale=0.7]
  \draw[thick,decorate,decoration={coil,aspect=0.4,segment length=4mm,amplitude=3mm}] (0,0) -- (2,0);
  \draw[thick,decorate,decoration={coil,aspect=0.4,segment length=1.5mm,amplitude=3mm}] (2,0) -- (3,0);
  \draw[thick,decorate,decoration={coil,aspect=0.4,segment length=4mm,amplitude=3mm}] (3,0) -- (6,0);
  \draw[<->,orange,thick] (1.9,0.6) -- (3.1,0.6);
  \node[orange] at (2.5,1.05) {\small compression};
  \draw[->,bleu,thick] (3.8,0.75) -- (5.8,0.75) node[above] {\small propagation};
\end{tikzpicture}
\end{minipage}
\end{cours}

\begin{cours}{Célérité et retard}
Un point $M'$ reproduit la perturbation du point $M$ avec un \textbf{retard} $\tau = t' - t$ ($t$ et $t'$ : instants d'arrivée en $M$ et en $M'$).
\[ \boxed{\,v = \frac{MM'}{\tau}\,} \qquad \text{soit} \qquad \tau = \frac{MM'}{v} \]
La \textbf{célérité} $v$ (vitesse de propagation, en $\ms$) dépend du \textbf{milieu}, pas de la forme ni de l'amplitude de la perturbation.
Son : environ $340~\ms$ dans l'air, $1\,500~\ms$ dans l'eau, $5\,000~\ms$ dans l'acier.
\end{cours}

\begin{cours}{Onde périodique : une double périodicité}
La source répète la perturbation à intervalles réguliers.
\begin{itemize}
  \item \textbf{Période $T$} (en s) : périodicité \textbf{temporelle}, la plus petite durée au bout de laquelle un point se retrouve dans le même état. \textbf{Fréquence} $f = \dfrac{1}{T}$ (en Hz).
  \item \textbf{Longueur d'onde $\lambda$} (en m) : périodicité \textbf{spatiale}, la plus petite distance entre deux points \textbf{en phase}.
  \item Pendant une période, l'onde avance d'une longueur d'onde :
  \[ \boxed{\,\lambda = v \times T = \frac{v}{f}\,} \]
  \item Deux points distants de $\lambda$, $2\lambda$, $3\lambda$\dots\ sont \textbf{en phase} (même état au même instant) ; distants de $\frac{\lambda}{2}$, $\frac{3\lambda}{2}$\dots\ ils sont en \textbf{opposition de phase} (l'un en haut quand l'autre est en bas).
\end{itemize}

\begin{minipage}{0.48\linewidth}
\centering
\begin{tikzpicture}[x=1.2cm,y=0.6cm]
  \draw[->] (0,0) -- (3.3,0) node[right] {$t$};
  \draw[->] (0,-1.3) -- (0,1.5) node[left] {$y$};
  \draw[thick,bleu,domain=0:3,samples=80] plot (\x,{sin(\x*240)});
  \draw[<->] (0.375,1.25) -- (1.875,1.25) node[midway,above] {$T$};
  \node[below] at (1.6,-1.3) {\small un point, en fonction du temps};
\end{tikzpicture}
\end{minipage}\hfill
\begin{minipage}{0.48\linewidth}
\centering
\begin{tikzpicture}[x=1.2cm,y=0.6cm]
  \draw[->] (0,0) -- (3.3,0) node[right] {$x$};
  \draw[->] (0,-1.3) -- (0,1.5) node[left] {$y$};
  \draw[thick,orange,domain=0:3,samples=80] plot (\x,{cos(\x*240)});
  \draw[<->] (0,1.25) -- (1.5,1.25) node[midway,above] {$\lambda$};
  \node[below] at (1.6,-1.3) {\small tout le milieu, à un instant};
\end{tikzpicture}
\end{minipage}
\end{cours}

\begin{piege}
\begin{itemize}
  \item $T$ se lit sur un graphique en fonction du \textbf{temps}, $\lambda$ sur un graphique en fonction de la \textbf{position}.
  \item Pour être précis, mesurer \textbf{plusieurs} périodes ou longueurs d'onde, puis diviser (5 intervalles entre 6 crêtes = $5\lambda$).
  \item La fréquence est fixée par la \textbf{source} : elle ne change pas quand l'onde change de milieu. La célérité, elle, change, donc $\lambda = \frac{v}{f}$ aussi.
  \item Convertir les unités : cm en m, ms en s, kHz en Hz, km/h en $\ms$.
  \item L'amplitude diminue en s'éloignant de la source (amortissement), mais la célérité ne change pas.
\end{itemize}
\end{piege}

\section{Exercices}

\exo{1}{Classer les ondes}
Pour chaque onde, dire si elle est transversale ou longitudinale, et si elle peut se propager dans le vide.\\
a) la ola dans un stade \quad b) une compression le long d'un ressort \quad c) le son \quad d) une vague à la surface de l'eau \quad e) une secousse le long d'une corde \quad f) la lumière du Soleil

\exo{1}{Calculs directs}
\begin{enumerate}
  \item Un son parcourt 1,02~km en 3,0~s. Calculer sa célérité.
  \item Pendant un orage, on entend le tonnerre 4,5~s après avoir vu l'éclair. À quelle distance est tombée la foudre ? ($v_{\text{son}} = 340~\ms$ ; la lumière arrive quasi instantanément.)
  \item Dans une cuve à ondes, la source vibre à $f = 20$~Hz et la longueur d'onde vaut 1,0~cm. Calculer la période et la célérité des ondes.
  \item Le «~la~» du diapason a une fréquence de 440~Hz. Calculer sa longueur d'onde dans l'air.
\end{enumerate}

\exo{2}{Une secousse le long d'une corde}
On crée en $S$ ($x = 0$), à $t = 0$, une secousse qui se propage le long d'une corde. On photographie la corde à deux instants.
\begin{center}
\begin{tikzpicture}[x=3.6cm,y=0.5cm]
  \foreach \yy/\tt/\xf in {2.6/{$t_1 = 0{,}10$ s}/0.60, 0/{$t_2 = 0{,}25$ s}/1.50} {
    \draw[thick] (0,\yy) -- (\xf-0.30,\yy) .. controls (\xf-0.2,\yy) and (\xf-0.2,\yy+1.2) .. (\xf-0.15,\yy+1.2) .. controls (\xf-0.1,\yy+1.2) and (\xf-0.1,\yy) .. (\xf,\yy) -- (2.2,\yy);
    \node[left] at (0,\yy) {\small\tt};
    \draw[dotted] (\xf,\yy) -- (\xf,-0.8);
  }
  \draw[->] (0,-0.8) -- (2.3,-0.8) node[right] {$x$ (m)};
  \foreach \i/\l in {0/0,0.3/{0,3},0.6/{0,6},0.9/{0,9},1.2/{1,2},1.5/{1,5},1.8/{1,8},2.1/{2,1}} \draw (\i,-0.7) -- (\i,-0.9) node[below] {\small\l};
\end{tikzpicture}
\end{center}
\begin{enumerate}
  \item L'onde est-elle transversale ou longitudinale ? Y a-t-il transport de matière ?
  \item À l'aide de la position du front de la secousse aux deux instants, calculer la célérité de l'onde.
  \item Un point $M$ est situé à 3,0~m de $S$. Avec quel retard reproduit-il le mouvement de $S$ ? À quel instant commence-t-il à bouger ?
  \item La secousse s'étend sur 0,30~m. Pendant combien de temps un point de la corde est-il en mouvement ?
\end{enumerate}

\exo{2}{La cuve à ondes}
Une soufflerie crée des vagues rectilignes à la surface de l'eau, à la fréquence $f = 25$~Hz. Vues du dessus, les crêtes apparaissent comme des lignes claires. On mesure \textbf{4,0~cm} entre la première et la sixième ligne claire. Trois petits flotteurs $A$, $B$ et $C$ sont posés sur l'eau, alignés dans la direction de propagation, avec $AB = 1{,}6$~cm et $AC = 2{,}0$~cm.
\begin{enumerate}
  \item Les flotteurs avancent-ils avec les vagues ? Que mettent-ils en évidence ?
  \item Calculer la période des vagues, puis leur longueur d'onde (attention : combien de longueurs d'onde entre la 1re et la 6e ligne ?).
  \item En déduire la célérité des vagues.
  \item $A$ et $B$ sont-ils en phase ou en opposition de phase ? Et $A$ et $C$ ? Justifier.
  \item Une crête passe par $A$ à l'instant $t_1$. Où se trouve-t-elle à l'instant $t_1 + T$ ?
\end{enumerate}

\exo{2}{Lire $T$ et $\lambda$ sur des graphiques}
Une onde sinusoïdale se propage le long d'une corde. Le graphique 1 donne l'élongation $y$ d'un point $M$ en fonction du temps ; le graphique 2 donne l'aspect de la corde à un instant donné.
\begin{center}
\begin{tikzpicture}[x=6cm,y=0.45cm]
  \draw[gray!30,very thin,xstep=0.1,ystep=1] (0,-2) grid (0.8,2);
  \draw[->] (0,0) -- (0.86,0) node[right] {$t$ (s)};
  \draw[->] (0,-2.2) -- (0,2.6) node[above] {$y$ (cm)};
  \foreach \i/\l in {0.2/{0,2},0.4/{0,4},0.6/{0,6},0.8/{0,8}} \node[below] at (\i,-2) {\small\l};
  \draw[thick,bleu,domain=0:0.8,samples=100] plot (\x,{2*sin(\x*900)});
  \node at (0.4,-3.3) {\small Graphique 1 : point $M$};
\end{tikzpicture}\hspace{8mm}
\begin{tikzpicture}[x=2cm,y=0.45cm]
  \draw[gray!30,very thin,xstep=0.3,ystep=1] (0,-2) grid (2.4,2);
  \draw[->] (0,0) -- (2.6,0) node[right] {$x$ (m)};
  \draw[->] (0,-2.2) -- (0,2.6) node[above] {$y$ (cm)};
  \foreach \i/\l in {0.6/{0,6},1.2/{1,2},1.8/{1,8},2.4/{2,4}} \node[below] at (\i,-2) {\small\l};
  \draw[thick,orange,domain=0:2.4,samples=100] plot (\x,{2*cos(\x*300)});
  \node at (1.2,-3.3) {\small Graphique 2 : la corde à un instant};
\end{tikzpicture}
\end{center}
\begin{enumerate}
  \item Sur quel graphique lit-on la période ? la longueur d'onde ? Lire leurs valeurs, ainsi que l'amplitude.
  \item Calculer la fréquence et la célérité de l'onde.
  \item Les points $N$ et $P$ sont situés respectivement à 3,6~m et 3,0~m de $M$, plus loin de la source. Calculer le retard de $N$ sur $M$. Comparer l'état vibratoire de $N$, puis de $P$, avec celui de $M$.
\end{enumerate}

\exo{2}{Sonar et échographie}
\begin{enumerate}
  \item Un sonar émet une salve d'ultrasons vers le fond de la mer et reçoit l'écho 0,40~s plus tard. Calculer la profondeur. ($v = 1\,500~\ms$ dans l'eau de mer.)
  \item Une sonde d'échographie émet des ultrasons de fréquence 5,0~MHz dans le corps, où leur célérité vaut $1\,540~\ms$. Calculer leur longueur d'onde en mm.
\end{enumerate}

\exo{3}{Localiser un séisme}
Un séisme produit deux types d'ondes : les ondes P (longitudinales, $v_P = 6{,}0~\text{km}\cdot\text{s}^{-1}$) et les ondes S (transversales, $v_S = 3{,}5~\text{km}\cdot\text{s}^{-1}$). Une station enregistre l'arrivée des ondes S \textbf{20~s} après celle des ondes P.
\begin{enumerate}
  \item Quelles ondes arrivent en premier ? Pourquoi ?
  \item On note $d$ la distance entre le foyer du séisme et la station. Exprimer les durées de trajet $\Delta t_P$ et $\Delta t_S$ en fonction de $d$.
  \item En déduire $d$.
  \item La station a enregistré les ondes P à 10~h~15~min~40~s. À quelle heure s'est produit le séisme ?
\end{enumerate}

\exo{3}{Mesurer la célérité des ultrasons}
Un émetteur envoie des ultrasons de fréquence $f = 40{,}0$~kHz vers deux récepteurs $R_1$ et $R_2$, dont on observe les signaux à l'oscilloscope. Au départ, $R_1$ et $R_2$ sont côte à côte et leurs signaux sont en phase. On éloigne $R_2$ dans l'axe : les signaux se retrouvent en phase 10 fois de suite, et on a alors déplacé $R_2$ de 8,5~cm.
\begin{enumerate}
  \item Quelle distance sépare deux positions successives de $R_2$ où les signaux sont en phase ?
  \item En déduire la longueur d'onde, puis la célérité des ultrasons.
  \item Pourquoi compter 10 décalages plutôt qu'un seul ?
  \item Les signaux sont-ils en phase ou en opposition de phase si l'on éloigne $R_2$ de 2,55~cm à partir de la position de départ ?
\end{enumerate}

\exo{3}{La houle arrive sur la plage}
Une houle de période $T = 8{,}0$~s se propage vers la côte. Loin de la plage, sa célérité vaut $12{,}5~\ms$. Près de la plage, elle n'est plus que de $5{,}0~\ms$.
En eau peu profonde, la célérité de la houle dépend de la profondeur $h$ : $v = \sqrt{g\,h}$ avec $g = 9{,}81~\text{m}\cdot\text{s}^{-2}$.
\begin{enumerate}
  \item Calculer la fréquence et la longueur d'onde de la houle loin de la plage.
  \item La période change-t-elle à l'approche de la plage ? Calculer la nouvelle longueur d'onde. Que constate-t-on ?
  \item Calculer la profondeur aux deux endroits.
  \item Un surfeur immobile voit passer combien de crêtes par minute ?
\end{enumerate}

\newpage
\section{Résultats}
\textit{Seulement les résultats, pour vérifier ses réponses : la rédaction reste à faire.}

\res{Classer les ondes}
Transversales : a, d, e. Longitudinales : b, c. Aucune ne se propage dans le vide, sauf f : la lumière est une onde électromagnétique (pas une onde mécanique).

\res{Calculs directs}
1. $v = 340~\ms$. \qquad
2. $d = 340 \times 4{,}5 \approx 1{,}5 \times 10^3$~m. \qquad
3. $T = 0{,}050$~s ; $v = 0{,}20~\ms$. \qquad
4. $\lambda \approx 0{,}77$~m.

\res{Une secousse le long d'une corde}
1. Transversale ; pas de transport de matière. \qquad
2. $v = \dfrac{1{,}50 - 0{,}60}{0{,}25 - 0{,}10} = 6{,}0~\ms$. \\
3. $\tau = 0{,}50$~s : $M$ commence à bouger à $t = 0{,}50$~s. \qquad
4. $0{,}30 / 6{,}0 = 0{,}050$~s.

\res{La cuve à ondes}
1. Non : ils montent et descendent sur place (pas de transport de matière). \qquad
2. $T = 0{,}040$~s ; 5 longueurs d'onde, donc $\lambda = 0{,}80$~cm. \\
3. $v = 0{,}20~\ms$. \qquad
4. $AB = 2\lambda$ : en phase ; $AC = 2{,}5\,\lambda$ : opposition de phase. \qquad
5. À 0,80~cm plus loin, sur la crête suivante.

\res{Lire $T$ et $\lambda$ sur des graphiques}
1. $T = 0{,}40$~s (graphique 1) ; $\lambda = 1{,}2$~m (graphique 2) ; amplitude 2~cm. \qquad
2. $f = 2{,}5$~Hz ; $v = 3{,}0~\ms$. \\
3. Retard de $N$ : 1,2~s $= 3T$. $MN = 3\lambda$ : $N$ en phase avec $M$ ; $MP = 2{,}5\,\lambda$ : $P$ en opposition de phase.

\res{Sonar et échographie}
1. Profondeur : $\dfrac{1\,500 \times 0{,}40}{2} = 300$~m. \qquad
2. $\lambda \approx 3{,}1 \times 10^{-4}$~m $= 0{,}31$~mm.

\res{Localiser un séisme}
1. Les ondes P (plus rapides). \qquad
2. $\Delta t_P = \dfrac{d}{v_P}$ et $\Delta t_S = \dfrac{d}{v_S}$. \qquad
3. $\dfrac{d}{3{,}5} - \dfrac{d}{6{,}0} = 20$ donne $d = 168$~km $\approx 1{,}7 \times 10^2$~km. \\
4. $\Delta t_P = 28$~s : le séisme s'est produit à 10~h~15~min~12~s.

\res{Mesurer la célérité des ultrasons}
1. Une longueur d'onde. \qquad
2. $\lambda = 0{,}85$~cm ; $v = 340~\ms$. \qquad
3. Pour diviser l'erreur de mesure par 10. \qquad
4. $2{,}55~\text{cm} = 3\lambda$ : en phase.

\res{La houle arrive sur la plage}
1. $f = 0{,}125$~Hz ; $\lambda = 100$~m. \qquad
2. Non ($T$ est fixée par la source) ; $\lambda = 40$~m : les crêtes se rapprochent. \\
3. $h \approx 16$~m loin de la plage ; $h \approx 2{,}5$~m près de la plage. \qquad
4. $60 / 8{,}0 = 7{,}5$ crêtes par minute.

\end{document}
